\section{MIFX Package}

A MIFX file is a self-contained archive containing machining process data and associated resources.

The archive uses the file extension:

\begin{verbatim}
.mifx
\end{verbatim}

A MIFX package is a standard ZIP archive containing a structured filesystem layout.

\subsection{Entry Point}

The root of every MIFX package must contain a file named:

\begin{verbatim}
job.json
\end{verbatim}

The \texttt{job.json} file acts as the entry point of the package and defines the primary structure of the machining job.

Software reading a MIFX package must begin interpretation from this file.

\subsection{Filesystem Organization}

The contents of a MIFX package are organized into directories representing different aspects of the machining process.

A typical package structure is shown below:

\begin{verbatim}
job.mifx
|
|-- job.json
|
|-- setups/
|   `-- set-1/
|       |-- setup.json
|       `-- operations/
|           |-- op-1.json
|           `-- op-2.json
|
|-- tools/
|-- geometry/
|-- in_process/
`-- addons/
\end{verbatim}

Each directory contains objects and resources associated with the machining job.

The structure is designed to remain human-readable while allowing software systems to efficiently navigate the package.

\subsection{Portability}

A MIFX package is intended to be portable between systems without modification.

All files required to interpret the machining job should be included within the archive.

External references are discouraged to ensure long-term integrity and reproducibility of the data.